% page 248 du fac-simile (image p-247 du scan)
% bloc 157.5 x 242.0 mm ; marge gauche 8.0 mm
\pgc{-12.700mm}{0.000mm}{
\l{\cel{3}240.}
\l{}
\l{}
\l{\sou{interpelar}. (\sou{trans., pri, pro}).I.Paroleskar (vers ulu, ad}
\l{\cel{3}ulu) por questionar lu pri ulo, o demandar ulo de lu.}
\l{\cel{3}- II. (\sou{yuro-cienco)}. Sumnar (ulu) justifikar su pri fak-}
\l{\cel{3}to. - III. (\sou{politiko}) (\sou{aludante parlamentano}). Demandar}
\l{\cel{3}a (ministro) ke ica justifikez, koram la parlamento,sua}
\l{\cel{3}agi o ne-agi. - DEFIRS.}
\l{}
\l{\sou{interpolar}. (\sou{trans.)}\cel{2}I.Insertar aden texto (frazo qua ne}
\l{\cel{3}esis parto di olu). - II. (\sou{fiziko}). Insertar aden observo-}
\l{\cel{3}rezultaji terminin determinita da kalkulo. - III. (\sou{alge}-}
\l{\cel{3}\sou{bro}). Trovar formulo qua, pro satisfacar ula quanto de kazi}
\l{\cel{3}observita, povas remplasar provizore la +leyo (lego) di}
\l{\cel{3}la fenomeno. - DEFIRS.}
\l{}
\l{\sou{interpretar}. (\sou{trans.})\cel{2}I.Explikar to quo esas obskura e ambi-}
\l{\cel{3}gua en texto. - II. Dicar to quon signifikas (oraklo,}
\l{\cel{3}sonjajo).- III. (\sou{teatro)}.\cel{2}Reprezentar segun la intenci}
\l{\cel{3}da autoro. - IV. Rezumar, en la linguo quan komprenas la}
\l{\cel{3}askoltanti, to quon dicas persono en linguo stranjera}
\l{\cel{3}quan la askoltanti ne komprenas. DEFIS.}
\l{}
\l{\sou{interstico}. Tre mikra spaco vakua inter la parti di korpo.}
\l{\cel{3}- EFIS.}
\l{}
\l{\sou{intervalo}.\cel{2}I. Disto-quanto inter du loki. - II. Disto-quanto}
\l{\cel{3}inter du instanti, du epoki. - III. (\sou{muziko)}. Disto-quanto}
\l{\cel{3}qua separas son-uno de la sono qua sequas lu nemediate,}
\l{\cel{3}acensante de grava ad akuta, o decensante de akuta a grava.}
\l{\cel{3}- DEFIRS.}
\l{}
\l{\sou{interviuvar}. (\sou{trans.}) (\sou{aludante la mrnalisti)}. Vizitar perso-}
\l{\cel{3}nego famoza (aktuala), por questionar lu pri lua vivo-ma-}
\l{\cel{3}niero, lua agi, lua idei, lua intenci pri lua agi balda,}
\l{\cel{3}e c. - DEF.}
\l{}
\l{\sou{\textquotedbl{}intestat}\textquotedbl{}\cel{2}Qua ne facis sua testamento.}
\l{}
\l{\sou{intestino}. (\sou{anat.}) Vicero en la kavajo abdomina, tubo digest-}
\l{\cel{3}ala ek mukozi e muskuli, qua iras de la stomako til la}
\l{\cel{3}anuso. - EFIRS.}
\l{}
\l{\sou{intima}. I. Qua esas tote interna. - II. Qua atingas la esenco}
\l{\cel{3}di la kozi. - III. (\sou{aludante personi)}. Qua unionesas (ad}
\l{\cel{3}altra) strete. - DEFIRS.}
\l{}
\l{\sou{intonar}. (\sou{trans.)}\cel{1}Komencar kantar (ario) por donar la tono}
\l{\cel{3}a la personi cetera. - DEFIS.}
\l{}
\l{\sou{intoxikar}. (\sou{trans.})\cel{2}Igar (ulu) absorbar substanco venenoza.}
\l{\cel{3}DEFIRS.}
\l{}
\l{intracelula. (\sou{biol.}) Qua esas interne di celulo (opoze ad}
\l{\cel{3}\textquotedbl{}intercelula\textquotedbl{} : qua esas inter la celuli).}
}
